\begin{abstract}
Formal CA2 evaluates \textbf{PAKS}: predictive workload forecasting combined with adaptive Kubernetes scaling versus traditional reactive Horizontal Pod Autoscaler (HPA). The committed method is an LSTM (TensorFlow named in the resources table) trained on Google Cluster Trace / Alibaba Cluster Trace, an adaptive scaler that acts through the Kubernetes API, and a cluster hosted on AWS EC2 with S3 datasets and CloudWatch metrics, reporting MAE/RMSE, utilisation, response time, throughput, scaling latency, infrastructure cost, and SLA/availability.

That live AWS Kubernetes experiment is \textbf{not} complete in this artefact (no cluster was deployed this pass). What \emph{is} implemented is: (i)~a public Google Cluster Data \emph{v1} (2010, 7-hour, CC-BY) slice with documented provenance, used to train a vanilla LSTM (NumPy BPTT; TensorFlow fail-closed on CPython 3.14); (ii)~fail-closed loaders for GCT 2011/2019 and Alibaba; (iii)~a dry-run adaptive engine that emits \texttt{apps/v1} Scale PATCH bodies (\texttt{dryRun=All}) against an \texttt{autoscaling/v2} HPA metrics schema; (iv)~the formal metric columns with explicit TRACE vs SIMULATED tags. A prior NimbusGuard-framed MLP simulator on synthetic sine/spike loads is retained only as a labelled \textbf{proxy} and is not the binding baseline (binding baseline: Kubernetes HPA). No MAE/RMSE, cost, or SLA figure from the proxy CSV is claimed as a formal CA2 outcome.
\end{abstract}
